\documentclass[11pt]{article}

\usepackage[margin=1in]{geometry}
\usepackage{booktabs}
\usepackage{graphicx}
\usepackage{amsmath}
\usepackage{multirow}
\usepackage{caption}
\usepackage[colorlinks=true, allcolors=blue]{hyperref}

\title{Replicating Stambaugh (1999), \emph{Predictive Regressions}}
\author{Ashish Maheshwari \and Omar Anabtawi}
\date{FINM 32900, Full-Stack Quantitative Finance \\ August 2026}

\begin{document}
\maketitle

\begin{abstract}
\noindent
Stambaugh (1999) shows that the ordinary least squares slope in a return
predictability regression is biased upward in finite samples when the predictor
is persistent and its innovations are negatively correlated with return
innovations---the defining features of the dividend--price ratio. We rebuild the
data from CRSP, replicate the paper's Table 1, Table 2, and Figure 1 within a
stated tolerance, and extend every exhibit through 2024. Our persistence
estimates match the paper to the third decimal, our finite-sample $p$-value for
1927--1996 is 0.177 against the paper's 0.17, and all sixteen Bayesian posterior
cells match within roughly 0.02. Extending the sample, we find the problem the
paper identifies has grown more severe: in 1997--2024 the naive $p$-value is
0.015 while the finite-sample $p$-value is 0.149, a tenfold gap.
\end{abstract}

\section{The problem}

The standard test of return predictability regresses next period's excess
return on this period's dividend--price ratio,
\begin{equation}
r_{t+1} = \alpha + \beta x_t + u_{t+1},
\qquad
x_{t+1} = \theta + \rho x_t + v_{t+1},
\end{equation}
where $x_t$ is the dividend--price ratio and the innovations $(u_{t+1},
v_{t+1})$ are jointly distributed with covariance matrix $\Sigma$. OLS is
consistent here, but it is not unbiased: strict exogeneity fails, because the
return innovation $u_{t+1}$ is correlated with the yield innovation $v_{t+1}$,
which in turn moves the \emph{future} regressor.

The bias arrives in two steps. First, the least squares estimator of an AR(1)
coefficient is biased downward by approximately $(1+3\rho)/T$ (Kendall, 1954);
with $\rho$ near one this is substantial. Second, that error transmits into the
predictive slope with its sign reversed,
\begin{equation}
E[\hat\beta - \beta]
= \frac{\sigma_{uv}}{\sigma_{vv}}\, E[\hat\rho - \rho]
\approx -\frac{\sigma_{uv}}{\sigma_{vv}} \cdot \frac{1+3\rho}{T}.
\label{eq:bias}
\end{equation}
Because a price increase raises the return and lowers the yield in the same
month, $\sigma_{uv} < 0$, and the slope bias is therefore \emph{positive}. Both
ingredients are necessary: remove the persistence or the innovation correlation
and the bias disappears. In our data both are present in abundance---persistence
of 0.99 and an innovation correlation of $-0.95$.

\section{Data}

Two WRDS pulls supply everything.

\paragraph{Market index.} We use \texttt{crsp.msi}, CRSP's pre-aggregated
monthly value-weighted market index, rather than the security-level stock file.
CRSP has already performed the capitalization weighting, so we consume a
finished series instead of reconstructing it from thousands of monthly
security-level observations, which would require lagged market-capitalization
weights and delisting-return handling.

\paragraph{Dividends without a dividend database.} The index file reports the
value-weighted return both including dividends (\texttt{vwretd}) and excluding
them (\texttt{vwretx}). Their difference is the month's dividend as a fraction
of the prior month's price. Compounding the price-only return yields a price
level $P_t$, so the monthly dividend in index units is
$D_t = (\text{vwretd}_t - \text{vwretx}_t) P_{t-1}$; summing twelve trailing
months gives the annual dividend, and the predictor is $D_{12,t}/P_t$.

\paragraph{Risk-free rate.} From \texttt{ff.factors\_monthly}. The dependent
variable is the continuously compounded excess return,
$\log(1+R_m) - \log(1+R_f)$, matching the paper's Table 1 specification.

\paragraph{Sample.} The merged panel spans June 1927 to December 2024
(1{,}171 monthly observations). It begins in mid-1927 because the risk-free
series starts in July 1926 and the trailing-twelve-month dividend window
consumes a further eleven months, leaving our replication samples about six
observations shorter than the paper's.

\paragraph{Index universe.} Stambaugh uses a NYSE-only value-weighted index.
Our WRDS instance provides no pre-built NYSE-only monthly index carrying both
\texttt{vwretd} and \texttt{vwretx}, so we use the CRSP total-market
value-weighted index. Because both are value-weighted, the same
large-capitalization firms---predominantly NYSE-listed---dominate each, and for
the first half of the sample the universes effectively coincide, since AMEX
enters CRSP in 1962 and NASDAQ at the end of 1972. A NYSE-only index rebuilt
from the stock file is reported separately as a robustness check.

\section{Replication results}

\input{../_output/table_01.tex}

Part C of Table~\ref{tab:table1} reproduces the paper's parameter estimates
closely: our persistence estimates of 0.973, 0.946, 0.981, and 0.990 compare to
the paper's 0.972, 0.948, 0.980, and 0.987, and the innovation covariances match
in sign, magnitude, and cross-sample pattern. Part A, simulated at those
parameters under the null $\beta = 0$, reproduces the paper's finite-sample bias
(0.074 against 0.07 in the full sample), standard deviation (0.155 against
0.16), skewness (0.74 against 0.71), kurtosis (3.90 against 3.84), and
finite-sample $p$-value (0.177 against 0.17). The contrast between Parts A and B
is the paper's central point: the standard regression model asserts an unbiased,
symmetric, normal estimator and delivers a $p$-value of 0.06, while the actual
finite-sample distribution is shifted, right-skewed, fat-tailed, and yields
0.18.

\input{../_output/table_02.tex}

Table~\ref{tab:table2} reports Bayesian posterior moments under the paper's four
specifications. Parts A and B use the likelihood conditioned on $x_0$; parts C
and D use the exact likelihood, which treats $x_0$ as a draw from the
predictor's stationary distribution, $x_0 \sim
N\!\left(\theta/(1-\rho), \sigma_v^2/(1-\rho^2)\right)$.

Under a flat prior with the conditional likelihood (Part A) the posterior is
conjugate and centers on the OLS estimate; the posterior mean of 0.21 and
$P(\beta \le 0) = 0.06$ match the paper exactly and provide an internal check
that the sampler is correct. Part B imposes stationarity, $\rho \in (-1,1)$, by
rejection. The restriction is nearly non-binding in the long samples but
discards roughly one draw in five in 1977--1996, where it raises the posterior
mean and halves $P(\beta \le 0)$ from 0.28 to 0.11 (the paper reports 0.26 to
0.13). The prior matters precisely where the data cannot identify persistence on
their own.

The exact likelihood in Parts C and D is nonlinear in $\rho$ and therefore not
conjugate, so these two specifications are sampled by random-walk
Metropolis-Hastings. The chain moves in an unconstrained parameterization---the
regression coefficients together with $\log\sigma_u$, $\log\sigma_v$, and
$\mathrm{atanh}\,\mathrm{corr}(u,v)$---so every proposal yields a valid
covariance matrix, with the corresponding Jacobian included in the log target.
The proposal covariance is estimated from a short pilot chain and scaled by the
usual $2.4/\sqrt{d}$ rule, because $\beta$ and $\rho$ are strongly correlated in
this posterior and a diagonal random walk mixes poorly along that ridge. We
validate the sampler by running it on the \emph{conditional} likelihood with a
flat prior, which is exactly the conjugate setting of Part A: it recovers a
posterior mean of 0.209, standard deviation 0.135, and $P(\beta \le 0) = 0.059$
against the conjugate values 0.21, 0.14, and 0.06.

The two adjustments push in opposite directions. The exact likelihood is
informative about $\rho$ and strengthens the evidence for predictability: for
1927--1996 the posterior mean rises to 0.23 and $P(\beta \le 0)$ falls to 0.04,
against the paper's 0.23 and 0.05. Part D's alternative prior, $p(b,\Sigma)
\propto (1-\rho^2)^{-1}\sigma_v^2|\Sigma|^{-5/2}$, places more weight near
$\rho = 1$ and pulls back: the mean falls to 0.17 and $P(\beta \le 0)$ rises to
0.13, against the paper's 0.19 and 0.10. All sixteen posterior cells agree with
the paper to within about 0.03. Chains were run to effective sample sizes above
200, which supports the means and tail probabilities to the two decimals
reported; the higher moments should be read as indicative only.

\begin{figure}[htbp]
\centering
\includegraphics[width=\textwidth]{../_output/figure_01.png}
\caption{Estimates of $\beta$ and $\rho$ by subsample and method. In every panel
the bias-adjusted estimate lies below and to the right of the OLS estimate:
correcting the finite-sample bias simultaneously lowers the slope and raises the
persistence, because both are manifestations of the same problem. The
1977--1996 panel reproduces the paper's most extreme result---the bias-corrected
persistence exceeds unity and the adjusted slope turns negative---which is the
correction reporting that this short, near-unit-root sample cannot identify the
parameters, and which is precisely the region the Bayesian stationarity
restriction excludes by construction.}
\label{fig:figure1}
\end{figure}

Figure~\ref{fig:figure1} reproduces the geometry of the paper's Figure 1. Our
bias-adjusted points correspond to the paper's ``F'' points and match closely:
$(\rho, \beta) = (0.978, 0.14)$ against the paper's $(0.976, 0.145)$ for
1927--1996, and $(1.007, -0.21)$ against $(1.007, -0.22)$ for 1977--1996.

\section{Extension to 2024}

\input{../_output/table_01_updated.tex}

\input{../_output/table_02_updated.tex}

Three results stand out. First, the bias behaves as theory predicts: in the full
sample $T$ rises from 834 to 1{,}170 months and the simulated bias falls from
0.074 to 0.057, tracking the $1/T$ rate in equation~(\ref{eq:bias}).

Second, modern slopes are not comparable to the paper's in raw units. The
1997--2024 slope of 1.43 is roughly seven times the paper's full-sample 0.21,
but the variance of the predictor's innovations falls over the same period from
$0.112 \times 10^{-4}$ to $0.010 \times 10^{-4}$. As firms substituted
repurchases for dividends the dividend yield declined to roughly 1.5\% and
ceased to vary. Since a slope measures return per unit of yield, a threefold
contraction in the predictor's standard deviation mechanically inflates the
coefficient. A standardized slope, $\hat\beta \times \mathrm{sd}(x)$, would be
the appropriate cross-era comparison.

Third, the paper's warning binds more tightly today than in 1999. In 1997--2024
the naive $p$-value is 0.015---conventionally significant---while the
finite-sample $p$-value is 0.149. The gap is tenfold, against roughly threefold
in the paper's own full sample, and the simulated bias of 0.567 is about 40\% of
the estimated slope. The Bayesian posterior for the same window places only 1\%
of its mass below zero, so the divergence between frameworks has widened as
well. We also note that the stationarity restriction has become nearly
irrelevant: the unrestricted sampler retains about 99\% of draws in every modern
window, against 80\% in 1977--1996, because longer samples pin down $\rho$
without help from the prior.

\section{Successes and challenges}

\paragraph{What replicated well.} The parameter estimates, the simulated
finite-sample moments, all sixteen Bayesian posterior cells across four
specifications, and the geometry of Figure 1 reproduce within a few hundredths,
including the paper's most extreme individual result.

\paragraph{Return definition.} Our initial slopes ran roughly 50\% high in the
pre-war subsamples. The cause was that we had used simple excess returns while
the paper's table note specifies continuously compounded returns; the difference
between the two is of the order of half the return variance, which is
substantial in the volatile 1927--1951 window and negligible later. Switching to
log returns moved the full-sample slope from 0.31 to 0.2100 against the paper's
0.21 and left $\rho$ and $\Sigma$ essentially unchanged, confirming the
diagnosis.

\paragraph{Vendor schema drift.} The course template's CRSP index puller targets
\texttt{crsp\_a\_indexes.msix}, the capitalization-decile file. The market index
we require resides in \texttt{crsp.msi} on our WRDS instance and uses
\texttt{date} rather than \texttt{caldt} as its date column. We resolved this by
enumerating libraries and describing tables through the \texttt{wrds} API before
writing the query rather than trusting documentation.

\paragraph{Sensitivity in the shortest sample.} Our simulated bias for
1977--1996 is 0.52 against the paper's 0.42. At $\hat\rho = 0.990$ with 239
observations the bias in equation~(\ref{eq:bias}) is extremely sensitive to
$\rho$, so a difference of 0.003 in the persistence estimate propagates
materially. We also condition the simulation on the first observed predictor
value, whereas the paper conditions on both sample endpoints.

\paragraph{A method that failed, and why.} Our first approach to Parts C and D
avoided building a sampler altogether: since the exact posterior differs from
the conditional one by a single observation's contribution, we reweighted the
Part B draws by the stationary density of $x_0$, using importance sampling. In
the long, lower-persistence samples this worked well, with effective sample
sizes above 90\% of nominal draws. In the post-war samples it failed. Part D's
prior carries a $(1-\rho^2)^{-1}$ factor that grows without bound as $\rho$
approaches one, so in the 1952--1996 sample the ten heaviest draws all had
$\rho > 0.999$ and the weighted posterior mean of $\rho$ was dragged from 0.981
to 0.989, displacing that specification's position in
Figure~\ref{fig:figure1}. The diagnosis is that the conditional posterior,
concentrated near $\rho = 0.98$, is a poor proposal for a target that places
substantial weight against the unit-root boundary; drawing more samples would
not have helped. We therefore replaced the reweighting with a
Metropolis-Hastings sampler that targets the exact posterior directly, which
recovers sensible persistence estimates (0.973 for Part C and 0.975 for Part D
in the full sample) and matches the paper across all sixteen cells. The
importance-sampling code remains in the repository, documented as the attempt
that motivated the switch.

\paragraph{Testing without credentials.} Our replication tests require the tidy
panel, which requires WRDS credentials that must never enter the repository.
The suite is therefore split: simulation-based tests of the bias mechanism run
in any environment, including continuous integration, while data-dependent
tests skip cleanly with an explanatory message when the panel is absent and run
locally once \texttt{doit} has built it.

\section{Conclusion}

The dividend--price ratio's apparent ability to predict returns is
substantially an artifact of estimation. In the paper's own sample the honest
$p$-value is three times the naive one; in the quarter-century since, it is ten
times. Whether any predictability survives depends on the inferential framework:
for 1927--1996 the finite-sample frequentist test cannot reject a zero slope at
0.18 while the Bayesian posterior places between 4\% and 10\% of its mass below
zero depending on the specification. That defensible frameworks disagree on
identical data is, we think, the paper's most durable lesson.

\section*{Reproducibility}

The project rebuilds end to end with \texttt{doit}, which pulls from WRDS,
constructs the tidy panel, regenerates every table and figure in this document,
executes the walkthrough notebook, and runs the test suite. Every statistic
reported here is generated by code; none is entered by hand.

\end{document}